\documentclass{article}
\usepackage[english]{babel}
\usepackage{booktabs}
\usepackage{multirow}
\usepackage{tabularx}

\begin{document}

% Case 1: Standard tabular with basic ccc columns and no borders
\begin{tabular}{ccc}
  1 & 2 & 3 \\
  4 & 5 & 6 \\
\end{tabular}

% Case 2: Standard tabular with lrc columns and single borders
\begin{tabular}{|l|r|c|}
  a & b & c \\
  d & e & f \\
\end{tabular}

% Case 3: Standard tabular with double vertical borders
\begin{tabular}{||c||c||}
  1 & 2 \\
  3 & 4 \\
\end{tabular}

% Case 4: Tabular with p{width} columns
\begin{tabular}{p{3cm}p{4.5cm}}
  Left text & Right text \\
  More left & More right \\
\end{tabular}

% Case 5: Tabular with mix of alignment and width
\begin{tabular}{|c|p{2cm}|l|}
  A & Width block & Left \\
  B & Content & Text \\
\end{tabular}

% Case 6: Tabular with hline at top, bottom, and middle
\begin{tabular}{cc}
  \hline
  1 & 2 \\
  \hline
  3 & 4 \\
  \hline
\end{tabular}

% Case 7: Tabular with double hline
\begin{tabular}{cc}
  \hline\hline
  A & B \\
  \hline\hline
  C & D \\
\end{tabular}

% Case 8: Tabular with cline{1-2}
\begin{tabular}{ccc}
  1 & 2 & 3 \\
  \cline{1-2}
  4 & 5 & 6 \\
\end{tabular}

% Case 9: Tabular with multiple clines on same row
\begin{tabular}{cccc}
  a & b & c & d \\
  \cline{1-1}\cline{3-3}
  e & f & g & h \\
\end{tabular}

% Case 10: Booktabs three-line table (toprule, midrule, bottomrule)
\begin{tabular}{cc}
  \toprule
  Header 1 & Header 2 \\
  \midrule
  Value 1 & Value 2 \\
  Value 3 & Value 4 \\
  \bottomrule
\end{tabular}

% Case 11: Booktabs three-line table with cmidrule{1-1} and cmidrule{2-2}
\begin{tabular}{cc}
  \toprule
  A & B \\
  \cmidrule{1-1}\cmidrule{2-2}
  C & D \\
  \bottomrule
\end{tabular}

% Case 12: Booktabs with thick/thin lines via custom trim
\begin{tabular}{cc}
  \toprule
  1 & 2 \\
  \cmidrule(r){1-1}
  3 & 4 \\
  \bottomrule
\end{tabular}

% Case 13: multicolumn spanning 3 columns with vertical borders
\begin{tabular}{|c|c|c|}
  \hline
  \multicolumn{3}{|c|}{Span Three Columns} \\
  \hline
  A & B & C \\
\end{tabular}

% Case 14: multicolumn spanning 2 columns with right alignment
\begin{tabular}{ccc}
  \multicolumn{2}{r}{Span Two Right} & C \\
  A & B & D \\
\end{tabular}

% Case 15: multirow spanning 2 rows
\begin{tabular}{cc}
  \multirow{2}{*}{Row Span} & A \\
                            & B \\
  C                         & D \\
\end{tabular}

% Case 16: multirow spanning 3 rows with custom width
\begin{tabular}{cc}
  \multirow{3}{2cm}{Fixed Width Row} & X \\
                                      & Y \\
                                      & Z \\
\end{tabular}

% Case 17: Multirow and multicolumn combined
\begin{tabular}{ccc}
  \multicolumn{2}{c}{\multirow{2}{*}{MultiRowCol}} & A \\
  \multicolumn{2}{c}{}                             & B \\
  C & D & E \\
\end{tabular}

% Case 18: Empty cell at the beginning of row
\begin{tabular}{ccc}
  & 2 & 3 \\
  4 & 5 & 6 \\
\end{tabular}

% Case 19: Empty cell at the end of row
\begin{tabular}{ccc}
  1 & 2 & \\
  4 & 5 & 6 \\
\end{tabular}

% Case 20: Multiple consecutive empty cells
\begin{tabular}{cccc}
  1 & & & 4 \\
  5 & 6 & 7 & 8 \\
\end{tabular}

% Case 21: Entirely empty row
\begin{tabular}{cc}
  & \\
  3 & 4 \\
\end{tabular}

% Case 22: Extra spaces and weird padding around & and \\
\begin{tabular}{  c  c  }
    1   &    2    \\
    3   &    4    \\
\end{tabular}

% Case 23: Tabularnewline instead of \\
\begin{tabular}{cc}
  1 & 2 \tabularnewline
  3 & 4 \tabularnewline
\end{tabular}

% Case 24: Nested tabular environments
\begin{tabular}{c}
  Outer cell \\
  \begin{tabular}{cc}
    Inner 1 & Inner 2 \\
  \end{tabular} \\
\end{tabular}

% Case 25: Table with math mode cells (inline $)
\begin{tabular}{cc}
  $x^2 + y^2$ & $\alpha \beta$ \\
  $\int_0^1 f(x) dx$ & $\sqrt{2}$ \\
\end{tabular}

% Case 26: Table with displaystyle block math
\begin{tabular}{c}
  $\displaystyle \sum_{i=1}^{n} i = \frac{n(n+1)}{2}$ \\
\end{tabular}

% Case 27: Table with text formatting inside cells
\begin{tabular}{cc}
  \textbf{Bold Text} & \textit{Italic Text} \\
  \underline{Underlined} & Plain \\
\end{tabular}

% Case 28: Table with special LaTeX character symbols
\begin{tabular}{cc}
  \% & \& \\
  \_ & \# \\
\end{tabular}

% Case 29: Extremely long text in cells
\begin{tabular}{p{5cm}}
  This is an extremely long sentence to test how natural word wrapping is handled inside a table cell when exporting or importing documents. \\
\end{tabular}

% Case 30: Very many columns (12 columns)
\begin{tabular}{cccccccccccc}
  1 & 2 & 3 & 4 & 5 & 6 & 7 & 8 & 9 & 10 & 11 & 12 \\
\end{tabular}

% Case 31: Very many rows (10 rows)
\begin{tabular}{c}
  Row 1 \\
  Row 2 \\
  Row 3 \\
  Row 4 \\
  Row 5 \\
  Row 6 \\
  Row 7 \\
  Row 8 \\
  Row 9 & Row 10 \\
\end{tabular}

% Case 32: Table wrapped in standard table float environment
\begin{table}[h]
  \begin{tabular}{cc}
    1 & 2 \\
    3 & 4 \\
  \end{tabular}
\end{table}

% Case 33: Table float with centering
\begin{table}[htbp]
  \centering
  \begin{tabular}{cc}
    A & B \\
  \end{tabular}
\end{table}

% Case 34: Table float with caption and label
\begin{table}
  \centering
  \begin{tabular}{cc}
    1 & 2 \\
  \end{tabular}
  \caption{Test Caption}
  \label{tab:test_label}
\end{table}

% Case 35: Table with vertical alignments t and b
\begin{tabular}[t]{cc}
  1 & 2 \\
\end{tabular}

% Case 36: Table with [b] vertical alignment
\begin{tabular}[b]{cc}
  3 & 4 \\
\end{tabular}

% Case 37: Multiple tabular environments inside one table float
\begin{table}
  \begin{tabular}{c}
    Table A \\
  \end{tabular}
  \begin{tabular}{c}
    Table B \\
  \end{tabular}
\end{table}

% Case 38: Comments inside tabular environment
\begin{tabular}{cc}
  % This is a comment row
  1 & 2 \\
  % Another comment
  3 & 4 % inline comment
\end{tabular}

% Case 39: Single cell table
\begin{tabular}{c}
  Solo \\
\end{tabular}

% Case 40: Table with no data cells but hlines
\begin{tabular}{c}
  \hline
  \hline
\end{tabular}

% Case 41: Multirow inside a tabular with top and bottom borders (test three-line support with multirow)
\begin{tabular}{cc}
  \toprule
  \multirow{2}{*}{DualHeader} & Header B \\
                              & Header C \\
  \midrule
  Value 1 & Value 2 \\
  \bottomrule
\end{tabular}

% Case 42: Tabular with multiple independent multirows on the same row
\begin{tabular}{ccc}
  \multirow{2}{*}{Multi 1} & \multirow{2}{*}{Multi 2} & Simple \\
                           &                          & Simple 2 \\
\end{tabular}

% Case 43: Tabular with multirow spanning 4 rows
\begin{tabular}{cc}
  \multirow{4}{*}{Span Four Rows} & R1 \\
                                  & R2 \\
                                  & R3 \\
                                  & R4 \\
\end{tabular}

% Case 44: Multirow inside multicolumn with custom width and border
\begin{tabular}{|ccc|}
  \hline
  \multicolumn{2}{|c|}{\multirow{2}{*}{Combined MultiRowCol Width}} & Col 3 \\
  \multicolumn{2}{|c|}{}                                            & Col 3 Row 2 \\
  \hline
\end{tabular}

% Case 45: Extreme multirow and multicolumn nested complex table
\begin{tabular}{|c|c|c|c|}
  \hline
  \multicolumn{3}{|c|}{\multirow{3}{*}{Extreme Nested Combined}} & Header 4 \\
  \multicolumn{3}{|c|}{}                                         & Sub 4-1 \\
  \multicolumn{3}{|c|}{}                                         & Sub 4-2 \\
  \hline
  Val 1 & Val 2 & Val 3 & Val 4 \\
  \hline
\end{tabular}

\end{document}
